\subsubsection*{3.1.6. 推广。}
前面的讨论可推广到如下情形：$f:X\to D$ 为一个态射，使 $0\in X$ 是\underline{在} $X_0$ 中的孤立非光滑点。缩小 $X$ 和 $D$ 后，$f$ 的非光滑点集在 $D$ 上有限。消灭圈流形 $V$ 具有有限个临界点构成的集合 $\Sigma$；$V-\Sigma$ 是带边界的可微流形（边界紧，但流形不紧）。单值变换 $T$ 在边界上为恒等，并保持 $\Sigma$。最后，命题 3.1.5 仍成立：保持纤维的同胚 $\Psi$ 在临界点集之外是微分同胚，且将临界点集映至 $W$ 中来自 $\Sigma$ 的部分。

\subsubsection*{3.1.7.}
命题 3.1.5 允许使用 1.4 的计算。得到：

\begin{enumerate}[label=(\Alph*)]
\item 消灭圈群 $(R^i\Phi(\Lambda))_0$ 是消灭圈流形 $V$ 以 $\Lambda$ 为系数的约化上同调群。
\item 群 $H^i_{\{0\}}(X_0,R\Psi(\Lambda))$ 是 $V^\circ$ 以 $\Lambda$ 为系数的具有固有支集的上同调群。
\item 变差
\[
  \operatorname{Var}:\widetilde H^i(V,\Lambda)
  \longrightarrow H_c^i(V^\circ,\Lambda)
  =H^i(V\bmod\partial V,\Lambda)
  \]
是由 $V$ 的自同构 $T$ 定义的经典变差，此自同构在 $\partial V$ 上为恒等。
\end{enumerate}

\subsubsection*{3.1.8.}
断言 (A)、(B) 和 (C) 实际上比命题 3.1.5 更初等。下面给出直接的证明与构造，从而无须验证 (A) 和 (B) 中的同构不依赖于 3.1.5 构造过程中的选择。

设 $d\leq r$ 及 $t\leq\tau(d)$，并设 $B$ 是开盒形区域
\[
B=\{x\mid\delta(x)<d\ \text{且}\ |f(x)|^2<t\}.
\]
设 $\widetilde D^*$ 是 $D^*$ 的万有覆盖，考虑将 (1.2.7.2) 限制到 $D_t$ 所得的图
\begin{equation}\tag{3.1.8.1}
\begin{tikzcd}[column sep=large,row sep=large]
\{0\} \arrow[r,hook,"\bar\imath"] \arrow[d,equal]
  & \widetilde D_t \arrow[d,"p"]
  & \widetilde D_t^* \arrow[l,hook',"\bar\jmath"'] \arrow[d,"p'"] \\
\{0\} \arrow[r,hook,"i"]
  & D_t
  & D_t^* \arrow[l,hook',"j"']
\end{tikzcd}
\end{equation}
设
\begin{equation}\tag{3.1.8.2}
\begin{tikzcd}[column sep=large,row sep=large]
B_0 \arrow[r,hook,"\bar\imath"] \arrow[d,equal]
  & \overline B \arrow[d,"p"]
  & \overline B^* \arrow[l,hook',"\bar\jmath"'] \arrow[d,"p'"] \\
B_0 \arrow[r,hook,"i"]
  & B
  & B^* \arrow[l,hook',"j"']
\end{tikzcd}
\end{equation}
是 $f:B\to D$ 在 (3.1.8.1) 上的逆像。仍用 $f$ 表示从 (3.1.8.2) 到 (3.1.8.1) 的投影。

\medskip
\noindent\textup{(A)}
按定义，$R^i\Psi_\eta(\Lambda)_0$ 是当 $U$ 遍历 $X$ 中 $0$ 的各邻域时，上同调群
\[
H^i(U\times_D\widetilde D^*,\Lambda).
\]
的正极限。令 $U$ 遍历由上述盒形区域 $B$ 构成的 $0$ 的基本邻域系。纤维积 $U\times_D\widetilde D^*=\overline B^*$ 是可缩空间 $\widetilde D_t^*$ 上以 $V^\circ$ 为纤维的丛。因此群 $H^i(\overline B^*,\Lambda)$ 的正向系统为常值，取值
\[
H^i(V^\circ,\Lambda)=H^i(V,\Lambda).
\]
所得同构
\begin{equation}\tag{3.1.8.3}
R^i\Psi_\eta(\Lambda)_0\simeq H^i(V,\Lambda)
\end{equation}
可进一步写成如下复合同构（参见下文 (B)）：
\begin{equation}\tag{3.1.8.4}
\begin{tikzcd}[column sep=small,row sep=large,cells={nodes={font=\small}}]
R\Psi_\eta(\Lambda)_0
  & R\Gamma(B_0,R\Psi_\eta(\Lambda)) \arrow[l,"\sim"']
  & (Rf_*R\bar\jmath_*\Lambda)_0 \arrow[l,"\sim"']
  & R\Gamma(\widetilde D_t,Rf_*R\bar\jmath_*\Lambda)
      \arrow[l,"\sim"'] \arrow[d,equal] \\
{}
  & R\Gamma(V,\Lambda) \arrow[r,"\sim"]
  & R\Gamma(V^\circ,\Lambda)
  & R\Gamma(\overline B^*,\Lambda) \arrow[l,"\sim"']
\end{tikzcd}
\end{equation}
复形 $R\Phi(\Lambda)_0$ 由区别三角
\[
\Lambda\longrightarrow(R\Psi_\eta\Lambda)_0
\longrightarrow R\Phi(\Lambda)_0\longrightarrow{}.
\]
定义。因此，群 $R^i\Phi(\Lambda)_0$ 是 $V$ 的约化上同调群。

\medskip
\noindent\textup{(B)}
设 $\mathfrak F$ 是 $\overline B^*$ 的那些在 $\overline B$ 中的闭包在 $\widetilde D$ 上固有的闭子集构成的族。下图中的箭头都是同构：
\begin{equation}\tag{3.1.8.5}
\begin{tikzcd}[column sep=small,row sep=large,cells={nodes={font=\small}}]
R\Gamma_{\{0\}}(R\Psi_\eta(\Lambda))
  \arrow[r,"\sim","\text{\textcircled{1}}"']
  & R\Gamma_c(B_0,R\Psi_\eta(\Lambda))
  & (Rf_!R\bar\jmath_*\Lambda)_0
      \arrow[l,"\sim"',"\text{\textcircled{2}}"] \\
R\Gamma_c(V^\circ,\Lambda)
  & R\Gamma_{\mathfrak F}(\overline B^*,\Lambda)
      \arrow[l,"\sim"',"\text{\textcircled{5}}"]
      \arrow[r,equal,"\text{\textcircled{4}}"']
  & R\Gamma(\widetilde D_t,Rf_!R\bar\jmath_*\Lambda)
      \arrow[u,"\text{\textcircled{3}}","\wr"']
\end{tikzcd}
\end{equation}
事实上，\textcircled{1} 是同构，因为 $B_0$ 是锥体（3.1.2）且 $R^i\Psi_\eta(\Lambda)$ 在 $0$ 外为常值；\textcircled{2} 由固有基变换定理；\textcircled{3} 因为 $Rf_!R\bar\jmath_*\Lambda$ 在 $0$ 外为常值；\textcircled{4} 是复合函子的谱序列的导出范畴形式；\textcircled{5} 因为 $\widetilde D^*$ 可缩。

同构 (3.1.8.5) 给出所求的同构
\begin{equation}\tag{3.1.8.6}
R^i\Gamma_{\{0\}}(R\Psi_\eta(\Lambda))
\simeq H_c^i(V^\circ,\Lambda).
\end{equation}
将在 3.1.10 中验证 (C)。

